\chapter{Four Surfaces of Human-Agent Collaboration}
\label{ch:four-surfaces}

\epigraph{A surface is where judgment becomes legible enough to trust, or fast enough to skip. The design question is which one you are building.}{}

\section{Method: Three Searches, One Map}

The synthesis in this chapter draws on three independent literature surveys, run through separate retrieval skills --- a deep-research web client, a citation-integrity researcher, and a high-recall keyword engine --- each tasked without sight of the others' results, then reconciled. The methodology appendix (Appendix~\ref{ch:methodology}) records the protocol. Two findings recurred across all three passes and organise this chapter. First, every collaboration surface is converging on graduated, risk-tiered autonomy rather than a binary automated-versus-manual switch --- a working reinterpretation of the classical levels-of-automation literature as a per-request property, not a fixed system setting. Second, Herbert Clark's theory of conversational grounding and the CSCW awareness tradition recur as the connective tissue across surfaces that otherwise share no vocabulary.

The evidence base is uneven, and this chapter says so where it bites. Peer-reviewed HCI venues (CHI, CSCW, VIS, ISMAR) now carry empirical human-agent teaming studies with checkable citations, while several 2026 attributions in the trust-calibration and voice literatures came back as vendor or grey-literature sources with no verifiable venue. Where a claim rests on a foundational framework rather than a fresh empirical result, the framework is named and a URL-backed source that applies it is cited; specific unverifiable author-year pairs are omitted rather than dressed up. That asymmetry --- forum and desktop surfaces densely evidenced, mixed reality and voice comparatively thin --- is itself a finding, and it is reflected in the weighting below.

\section{Why Four Surfaces, Weighted}

An \meshterm{interaction surface} is the medium and modality through which a human and an agent exchange initiative: where the agent's activity becomes perceptible, where the human's judgment is entered, and where the two are grounded against a shared state. Four surfaces account for essentially all deployed human-agent collaboration in 2025--26: asynchronous text and forum, desktop immersive 3D, mixed reality, and voice. They are not equal in analytical weight, and this edition apportions attention deliberately (Table~\ref{tab:four-surfaces}).

\begin{definitionbox}[title={Graduated autonomy}]
\meshterm{Graduated autonomy} is the treatment of an agent's authority as a per-request property, assessed from the verifiability, reversibility, and blast radius of each proposed action, rather than a fixed setting for the whole system. Low-risk reversible actions run autonomously with logging; moderate-risk actions route to an asynchronous review queue; irreversible or high-impact actions require explicit human authorisation before execution. Every surface in this chapter is converging on this shape, which is why the design question ceases to be ``how autonomous is the agent'' and becomes ``on this surface, at this stake, who holds initiative for this action''.
\end{definitionbox}

\begin{table}[h]
\centering
\caption{The four surfaces, their collaborative function, and this edition's analytical weight}
\label{tab:four-surfaces}
\begin{tabularx}{\textwidth}{L{3.4cm} C{1.6cm} X L{3.0cm}}
\toprule
\textbf{Surface} & \textbf{Weight} & \textbf{Primary collaborative function} & \textbf{Governance role} \\
\midrule
Async text / forum & 40\% & Deliberate, legible, auditable judgment on parked agent actions & Decision locus \\
\addlinespace
Desktop immersive 3D & 30\% & Awareness and steering of work in flight & Observation plus intervention \\
\addlinespace
Mixed reality & 20\% & Co-present awareness at room and team scale & Presence, not decision \\
\addlinespace
Voice & 10\% & Low-friction command ingress & Ingress, not governance \\
\bottomrule
\end{tabularx}
\end{table}

The weighting follows from what each surface is good for. Judgment that must be legible, auditable, and deliberate belongs on the asynchronous surface, because text persists, carries its own provenance, and imposes no latency tax on the human --- the reviewer reads the parked action, the evidence, and the uncertainty at a time of their choosing and signs a durable record. Awareness and steering favour the visual surfaces: a running multi-agent workflow is a spatial, high-volume object that a 3D canvas renders glanceable in ways a text log cannot. Mixed reality extends that awareness to co-present, multi-user settings but adds no distinct decision semantics of its own. The conservative weight it receives here continues a longer-standing reading: the author's 2022 assessment judged social-first metaverse platforms the weakest of the remote-collaboration options and predicted that Meta's Horizon Worlds and Decentraland would fail to hold an audience, which happened within roughly a year \parencite{OHare2022Convergence}. Voice is the fastest way to start work and the worst way to govern it; it is command ingress, not a channel for the accountable decision.

\begin{figure}[h]
\centering
\includegraphics[width=0.92\textwidth]{nb-four-surfaces.jpg}
\caption{The four interaction surfaces and their analytical weights: judgment on the asynchronous surface, awareness on the visual surfaces, ingress by voice. Illustration in the house sketch style.}
\label{fig:four-surfaces}
\end{figure}

The remainder of this chapter takes each surface in turn under the same three headings: what published theory demands of it, what shipping products had made table-stakes by 2026, and the design questions still open. The order tracks the weighting, so the asynchronous surface --- where the accountable decision lives --- gets the fullest treatment, and the discussion of DreamLab's own canon stays light, held mostly for the component chapters that follow and the gap register in Chapter~\ref{ch:gap-register}.

\section{Asynchronous Text and Forum (40 Per Cent)}

\subsection{What theory demands}

The asynchronous surface inherits the richest theoretical apparatus. Horvitz's twelve principles of mixed-initiative interaction remain the base layer for coupling automated services to direct human control, weighing the utility of an agent's action against its uncertainty and the cost of interrupting the human \parencite{Horvitz1999Principles}. A 2025 scoping review finds those principles under-specify how much control today's agents actually claim, and that the field lacks consensus even on what a mixed-initiative system is --- Horvitz's synthesis was built to resolve the 1997 Shneiderman--Maes dispute between direct manipulation and interface-agent delegation, a dispute the agent era has reopened rather than closed \parencite{Monadjemi2025Scoping, ShneidermanMaes1997}. For coordination itself, Dourish and Bellotti's account of awareness passively conveyed through a shared workspace remains the seminal reference, extended operationally by Gutwin and Greenberg's descriptive framework; recent CSCW work re-applies both to agents treated as remote collaborators rather than tools, and finds current systems short on mutual awareness and shared accountability \parencite{DourishBellotti1992, GutwinGreenberg2002, HCIScoping2025}.

Delegation is where theory has moved fastest. The 2025--26 consensus reframes it as a dynamic, per-task assessment of verifiability and impact rather than a static permission, with an escalation gate that fires when the failure impact of an action exceeds a resilience threshold \parencite{Tomasev2026Delegation}. Trust research supplies the caution: calibrated trust, not maximal trust, predicts team success, and over-trust slides teams from cooperation into passive delegation and automation complacency \parencite{TrustingTeammates2025, CHAIT2025}. The systematic evidence on automation bias is blunt --- confidence in the system, not its accuracy, drives reliance, and more reliable systems produce less vigilant overseers \parencite{AutomationBias2012}.

Disclosure and provenance form the third leg. Work extending the W3C PROV standard captures agent prompts, decisions, and outputs for traceability \parencite{Souza2025ProvAgent}, while a cluster of 2026 papers maps the gap between decentralised-identity infrastructure and deployed agents and argues that the agentic web needs agents to identify themselves and their authorising principals as a condition of participation \parencite{AIIdentity2026, AuthenticatedDelegation2025, AgenticWebNorms2026}. This is the theoretical warrant for the identity model DreamLab's canon already assumes --- one signing key that is at once relay identity, HTTP-auth identity, access-control principal, and provenance author --- and the literature's open admission that no such model is yet standardised is what makes the canon's bet legible as a bet. Community-moderation research adds that multiple agents conversing as a group exert more social influence than a single agent, and that third-party moderation bots imposed from outside a community mismatch its norms --- Botender answers this by letting communities co-author agent behaviour through generated provocation scenarios before deployment \parencite{Song2025SocialGroups, ThirdPartyBots2024, Kuo2026Botender}.

Disclosure changes the social relationship as much as the audit trail. Scoping work on human-AI co-creation reports that once an agent's authorship is disclosed, human collaborators reframe it as a subordinate to manage rather than a colleague, adopting stances that range from teaching and orchestrating to resisting and complying; undisclosed authorship, by contrast, produces a measurable disadvantage for contributors who are wrongly assumed to be less proficient \parencite{HCIScoping2025}. For the asynchronous surface this is the practical case for making agent presence ambient and identity legible by construction: the awareness that Dourish and Bellotti describe as passively conveyed through a shared workspace is exactly what an agent's signed, provenance-carrying post supplies, provided the surface renders it rather than hiding it behind a human-looking byline.

\begin{evidencebox}[title={Escalation over approval, on the surface built for it}]
The 2025--26 delegation literature and this report's own deployment evidence point the same way: default to escalation, not blanket approval. Tomašev and colleagues formalise the escalation gate as firing only when an action's failure impact crosses a resilience threshold \parencite{Tomasev2026Delegation}, and the automation-bias record shows that gating every output trains overseers to stop looking \parencite{AutomationBias2012}. The asynchronous surface is where that discipline is cheapest to implement: a parked action, its evidence, and its uncertainty can wait for a considered human decision without stalling the agent's other work, which is precisely why the accountable decision belongs here and not on the faster surfaces.
\end{evidencebox}

\subsection{What shipped}

By 2026 the diff-and-approve gate is table-stakes across production agents. Claude Code's three permission modes --- ask-per-edit, auto-accept-edits, and a zero-side-effect plan mode --- are a product-level instantiation of the levels-of-automation taxonomy, though the vendors rarely cite it \parencite{ClaudeCodeDesktop2026}. GitHub's Agent HQ attempts to unify oversight across previously siloed agent sessions \parencite{AgentHQ2026}, and the async approval queue --- park the action, summarise it, continue other work, resume on approval --- has become the standard answer to the latency tax that synchronous stop-and-wait approval would impose at scale.

Governance of agent \emph{identity}, by contrast, remains reactive and crisis-driven. The University of Zurich's undisclosed persuasion bots in r/changemyview and Wikipedia's 2026 block of an autonomous article-writing agent both show platforms improvising disclosure and ban policy after the fact, despite mixed-initiative disclosure principles having existed since 1999 \parencite{ZurichCMV2025, WikipediaAgentBan2026}. Regulation is closing the gap from the other side: the EU AI Act's high-risk oversight duties push human-in-the-loop from best practice to compliance obligation \parencite{EUAIAct2025}, and NIST's risk-management framework and ISO/IEC~42001 both now mandate explicit human authorisation for agentic actions \parencite{NISTAIRMF, ISO42001}.

\subsection{Open questions}

Approval fatigue is the surface's defining failure mode: operators enable auto-approve until conditional oversight degrades into routine bypass, the exact complacency the trust literature predicts. No ratified standard yet governs agent self-disclosure --- the identity work remains draft, and a 2025 index of thirty deployed agentic systems found disclosure practice inconsistent industry-wide \parencite{AgentIndex2025}. Multi-human and conflicting-decision resolution is essentially unstudied. This is the surface DreamLab's canon commits to most heavily: nostr-rust-forum is positioned as the sole fully-specified locus of human judgment, with only cryptographically authenticated humans permitted to sign the decision event. The 2026-07-03 closeout is candid that the relay's advertised NIP-42 authentication is in practice a pubkey whitelist with \texttt{auth\_required:false} --- the canon's trust claim for this surface currently overstates its own mechanism, and Chapter~\ref{ch:gap-register} logs it.

\emph{After the gap-close sprint.} The forum was the most fully built surface entering the sprint and closed most of its own remaining gaps. Agents now carry a disclosure badge at every author-render site, fed by a public active-only registry projection (\texttt{7157a92}); an ordinary member reaches a read-only governance view where the route had bounced all non-admins (\texttt{0b8a1c}); the \texttt{broker\_decisions} audit trail gained the read endpoint it lacked (\texttt{GET /api/governance/decisions}); graduated delegate, promote and precedent decisions now project through the relay with a risk tier attached (\texttt{6986276}); the client signs through a relay-confirmed \texttt{publish\_with\_ack} rather than reporting success before the OK frame; and a scaffolded, canon-gated supersession path lets an authorised signer mark a prior decision \texttt{Superseded} by a new signed reference (\texttt{35dbb1b}). What the surface still owes is full member participation in the decision itself and a cross-relay path, both of which stay open.

\section{Desktop Immersive 3D (30 Per Cent)}

\subsection{What theory demands}

Where the forum is built for deliberation, the desktop 3D surface is built for awareness of work in flight. AGDebugger is the clearest 2025 anchor: an overview visualisation clusters high-volume agent-to-agent message histories, and an edit-and-reset affordance lets a user pause, edit a prior agent message, and replay execution from that point, establishing interruptibility and checkpoint-resettability as core observability affordances \parencite{Epperson2025Debugging}. Trace-visualisation tools render per-node tool-call flows for glanceable failure attribution \parencite{XAgen2025}, hierarchical orchestration views and studies of interruptible long-horizon agents extend the same needs empirically \parencite{OrchVis2025, Interruptible2026}, and design work argues intent-communication modality must follow communicative purpose rather than defaulting to text \parencite{IntentComm2025}.

The organising dialectic is again Shneiderman versus Maes. Nielsen frames command interfaces as keeping the human's plan externalised and the system transparent, while intent-based delegation pushes the system into interpreting and subgoaling, with the classical supervisory-control failure mode --- the operator shifts from doer to monitor and loses signal-sensitivity \parencite{Nielsen2025Intent}. Endsley's three-level model of situation awareness, revisited for agent oversight, predicts that projective awareness decays fastest when agents update their world-models without notifying the human \parencite{SupervisionTeaming2025}. The design conclusion drawn across this literature is that a desktop surface should expose steering as an adjustable affordance --- from full human control through inner-loop autonomy to full delegation --- rather than a fixed binary approve-or-reject. Baldauf's delegation-interaction model renders that as a spectrum from action confirmation, close to direct manipulation, to supervisory co-execution, close to delegation, with risk-gated autonomy best preserving situation awareness without cognitive overload \parencite{Baldauf2026Delegation, ControlAbstraction2020}.

Preference data supports keeping the human in that loop. Workers asked how much autonomy they want per task cluster well short of full automation \parencite{HumanAgencyScale2025}, an organisational study found employees willing to cede only about a third of decision rights to AI, and interactive-plan systems critique the dominant preview-then-execute-then-review pattern as too thin for genuine co-work \parencite{Cocoa2025}. The recurring lesson is that legibility on this surface is not the same as verbosity: a 3D canvas earns its thirty per cent of weight only when it makes the shape of a running computation graspable at a glance, which a scrolling text trace cannot, and surrenders that advantage the moment it merely renders the log in three dimensions.

\subsection{What shipped}

Observability is mid-consolidation. OpenTelemetry has become the interoperability layer for agent traces, so a system can be instrumented once and routed to LangSmith, Phoenix, or Langfuse, and a formal AgentOps taxonomy specifies observability needs across user, deployer, and model-developer stakeholders \parencite{AgentOps2024}. Devin and Replit stream file-tree and console activity live rather than as post-hoc diffs --- a legibility-via-liveness pattern distinct from the forum's parked-diff review. Practitioner tools still outpace vendors for swarm-specific situation awareness: an open-source hooks dashboard that pipes every Claude Code tool call and subagent lifecycle event into a real-time WebSocket visualisation fills a gap the commercial platforms do not yet cover \parencite{DislerObservability2025}.

\subsection{Open questions}

Reliability science complicates the steering picture: agent reliability does not scale with capability --- larger models often show lower run-to-run consistency because they have more viable solution paths --- and how much that matters depends on whether the agent runs autonomously or augments a human, which bears directly on how tightly a steering UI should couple to task type \parencite{AgentReliability2026}. Token-cost-per-trace in multi-agent runs is five-to-ten times naive per-call estimates, a legibility problem specific to swarms. And projective situation awareness --- Endsley's Level~3 --- has no settled visual grammar. DreamLab's VisionClaw sits precisely here: its desktop client renders agent actions as transient beams and makes blocked-versus-approved states obvious through GPU semantic physics, but the canon assigns the surface no decision authority. Humans watching the graph observe the consequences of decisions made in the forum; they do not sign responses on the canvas, and the attractive-force ``gluon'' embodiment remains a documented no-op. Observation is shipped; steering is partial.

\emph{After the gap-close sprint.} The steering the surface owed is now mounted. VisionClaw's desktop opens an agent-detail panel on node-select, carries a broker case queue with an ambient cue when a judgment is pending, and shows a swarm-observability panel tagging failures to the fourteen-mode MAST taxonomy (\texttt{e0f582403}); an operator can submit, interrupt, and query a task, and where an agent is MCP-native and cannot be interrupted the surface discloses that rather than offering a dead button (\texttt{b6c1c43b5}). Agent nodes now carry a \texttt{did:nostr} minted at spawn and verified by a Schnorr check (\texttt{4a595cc8f}), and the live-connection dot that once reported healthy unconditionally now reads ground truth (\texttt{6f4eb1b0a}). The gluon attractive force stays deferred and the live beam remains a pending on-device canary, so observation is shipped and steering is now wired, with its last proof still to run on a live session.

\section{Mixed Reality (20 Per Cent)}

\subsection{What theory demands}

Mixed reality extends desktop awareness to co-present, room-scale, and multi-user settings. The evaluative constructs are co-presence (the sense of physically being together) and social presence (psychological connection through attention and emotional contagion), and the field's recent finding is that neither is strictly anthropomorphic: spatialised voice with situated movement sound raises co-presence even without a visible avatar, and inanimate objects that speak can drive higher engagement than humanoid guides in the right context \parencite{Schmid2024Objects}. A 2025 systematic review of conversational agents in augmented reality frames the trend as a move toward dynamic, context-aware embodiment along a presence-form by context-awareness grid \parencite{ARCASurvey2025}. The binding constraint the field now names is latency, not dialogue quality: below-conversational response times, not richer avatars, are what let free-flowing agent conversation feel co-present at all.

Proxemics and gaze carry structural load. Dynamic gaze-centring draws users physically closer to agents except in high-social-anxiety users, implying comfort-adaptive interpersonal distancing is a design requirement, and proxemic distance plus gaze can initiate and structure multi-party conversation in a mixed-initiative style \parencite{VirtualGaze2024, ProxemicsGaze2018, PreTouchGaze2023}. For asymmetric collaboration --- one immersed headset user, one desktop user --- an embodied mediator agent can point out spatial features to the VR user while giving data-rich summaries to the desktop partner, an explicit agent-as-information-bridge model; systematic reviews and layout-dimensionality studies map the design space of pairing a 2D operator with a 3D collaborator \parencite{AsymmetricSLR2025, LayoutDim2024}. Spatially-aware frameworks now have agents build a navmesh over real furniture for context-appropriate locomotion and distance \parencite{SpatialEmbodied2026}.

The author had worked with these constructs long before this literature existed. A fifteen-year telepresence research programme, run through the Octave mixed-reality laboratory at NICVE, produced a 2022 synthesis of the same gaze and proxemics work, and that synthesis had already fixed the numbers a credible mutual-gaze system must hold: an optical-axis-to-camera offset within 1.2 degrees horizontally and 1.7 degrees vertically, a fall of roughly a quarter in multiparty turn-passing efficiency once eye-gaze support is absent, and a correlation between synchronised gaze and task performance around 46 per cent higher \parencite{OHare2022Convergence}. It reached the anthropomorphism result from the business-collaboration side as well: across the B2B systems it surveyed, embodied telepresence robots and avatars matched but did not exceed a well-engineered face-on-screen video call in perceived personalness or interactivity, which is why this book ranks avatar fidelity below gaze and audio fidelity throughout the component chapters. The 2025--26 co-presence sources arrive at the same binding constraints by a different route, three years on.

\begin{figure}[h]
\centering
\includegraphics[width=0.72\textwidth]{b22-telepresence-venn.png}
\caption{The author's mixed-reality telepresence research map as published in 2022: communication theory, psychology, and technology intersecting on mutual gaze, proxemics, and presence effects \parencite{OHare2022Convergence}. Reproduced as published, including the original's struck-through ``immersive'' --- a revision mark the map carried in print. The constructs this chapter's mixed-reality section evaluates against are the ones this map fixed.}
\label{fig:telepresence-venn}
\end{figure}

\subsection{What shipped, and what has not}

This surface is the thinnest in deployed practice. Embodied-agent-as-a-service exists as a packaged SDK category --- avatar-rendering pipelines offering ``FaceTime with your AI agent'' \parencite{MatrixOne2026} --- but there is no dominant interaction pattern the way diff-approval dominates the forum. Asymmetric headset-plus-desktop collaboration is handled mainly through enterprise VR-meeting bridging rather than research-validated interaction patterns, and the academic work remains lab-scale. DreamLab's canon treats mixed reality identically to desktop but at room scale: the same binary protocol and \texttt{did:nostr} identity feed a native XR client validated at 250-plus concurrent users. That 250-user figure measures how many clients the transport sustains, a different property from whether those users experience one another as co-present. The author's own telepresence work found presence and interpersonal trust degrade well before a group reaches the dozens, and put the productive collaboration size at around six; the 2022 design rejected the social-metaverse model of hundreds of avatars on exactly that evidence \parencite{OHare2022Convergence}. So 250 concurrent connections is a transport result, and the co-presence quality it might seem to promise is what the register records as unrepresented in the codebase. Crucially, the canon gives XR no distinct governance semantics --- no text addresses XR-specific consent, latency-sensitive judgment, or what happens when two immersed users register conflicting decisions. Mixed reality is a presence surface here, and the absence of any XR governance model is one of the register's cleaner silences (Chapter~\ref{ch:gap-register}).

\emph{After the gap-close sprint.} The co-presence quality the register recorded as unrepresented in the codebase now has an implementation to validate. A Godot copresence layer (\texttt{485bb8f9f}) supplies a gaze abstraction that falls back from eye to head ray where the hardware lacks eye tracking, a three-resolver selection arbiter of ray, pinch, and gaze-dwell, a Hall-zone proxemics solver, an avatar state machine with mutual-gaze engagement, and a compact \texttt{0x44} presence codec (\texttt{57b32faee}); each agent renders with a DID badge that stays amber until a Schnorr check verifies it, and a HUD panel signs an intervention under NIP-98 without the key leaving the device (\texttt{0f3a1b60c}). It stands at \texttt{standalone} tier, validated in test but not on a physical headset, so on-device co-presence measurement is the pending item. The canon still writes no XR-specific governance for conflicting immersed decisions, which remains the surface's open silence.

\section{Voice (10 Per Cent)}

\subsection{What theory demands}

Voice research has moved from silence-based turn detection toward semantic turn-taking: conditioning floor-yielding on assigned conversational roles and reasoning over the audio stream cuts false-positive interruptions materially, and surveys trace the shift from hidden-Markov turn models to continuous duplex-audio architectures \parencite{AdaptiveTurnTaking2026, TurnTakingSurvey2025, ChronologicalThinking2025}.

Clark's grounding and repair theory is the explicit anchor for the surface's central weakness. LLM agents produce locally appropriate acknowledgments but fail to reliably reuse established common ground later in a dialogue, and issue clarifying repair requests far less often than humans do \parencite{GroundingRifts2025, RepairTaxonomy2025, ConvGrounding2024}. Multi-agent negotiation work extends grounding-failure analysis to settings where several agents share a channel \parencite{TalkCheap2026}, and benchmarks now test whether an agent in a three-or-more-party group can tell a turn-yield addressed to it from a pause inside a human-human sub-dialogue \parencite{SpeakStaySilent2026}. A trust-precision trade-off runs through the applied findings: users favour explicit command mode for high-stakes tasks and conversational mode for relational ones, implying agents need dynamic mode-switching rather than a single register.

\subsection{What shipped}

The half-cascade speech-to-speech stack --- speech-to-text, model, text-to-speech collapsed into one streaming connection --- is the 2026 reference architecture, with latency now the competitive axis at roughly 200--280 milliseconds rather than transcription accuracy \parencite{Deepgram2026}. Barge-in remains unreliable even in flagship APIs: an open bug thread documents the model over-triggering interruption on backchannel filler words (``hmm'', ``uh-huh'') that function as grounding acknowledgments, not turn-claims --- a direct failure to encode Clark's theory, and a case where even table-stakes voice UX is not solved \parencite{RealtimeBargeIn2026, RepairScoping2024}. Production multi-agent voice adds an orchestrator layer plus a tiered memory for coherent addressing across a session, but which agent the user is addressing when several share a channel has no established convention.

\subsection{Open questions and the governance line}

Multi-agent addressing, reliable barge-in, and grounding-reuse are all open. Latency binds this surface as it binds mixed reality, and the author's telepresence research anticipated that too: the 2022 synthesis named spatial-audio latency and directionality non-negotiable for co-present collaboration, a channel-fidelity finding distinct from the turn-taking latency at issue in command ingress, though it points to the same conclusion from a different surface \parencite{OHare2022Convergence}. The load-bearing point for this book is narrower: voice is ingress, not governance. DreamLab's canon makes this a deliberate least-authority commitment --- push-to-talk captures the currently-selected agent node and dispatches a scoped, signed action request to that agent's identity, never a channel for the human decision itself. The honest caveat, recorded in the closeout, is that this binding is weaker in code than in prose: the voice-to-selected-actor step lacks the file-and-line evidence that surrounds other shipped claims, and the contract it depends on is broken today because live agent nodes are keyed by \texttt{task\_id} rather than \texttt{did:nostr} and so cannot yet be addressed by the action-request protocol. Voice earns ten per cent of this edition's weight because it is where work starts, not where it is judged.

\emph{After the gap-close sprint.} The broken contract the closeout recorded is largely closed. A \texttt{VoiceIntentClient} now signs a \texttt{kind-31402} request under a NIP-98 mandate and routes it to the agent's own identity through agentbox's \texttt{/v1/voice-intent} producer, consuming the existing push-to-talk service rather than superseding it (\texttt{f6e1d58f9}), and a confidence gate asks for clarification instead of acting on a low-confidence transcript (\texttt{774ffa05e}). The addressing rests on the \texttt{did:nostr} keying the sprint carried end to end (\texttt{4a595cc8f}) rather than the \texttt{task\_id} that had made an addressed actor unaddressable, so voice is now a scoped, signed ingress to a named agent. The one item left is on-device proof of the end-to-end round-trip, seeded as a canary.

\section{Two Threads, and the Register Ahead}

\begin{keyinsight}[title={What Every Surface Shares}]
Across all four surfaces, two convergences hold. Autonomy is becoming graduated and per-request --- escalation gates in forums, risk-gated steering on the desktop, comfort-adaptive proxemics in mixed reality, command-versus-conversation switching in voice --- a 2020s reinterpretation of the levels-of-automation model as a dynamic property rather than a static setting. And Clark's grounding, with the CSCW awareness tradition, is the connective tissue: forum agents need explicit common-ground status updates, desktop debuggers visualise shared-mental-model breakdown, mixed-reality mediators bridge grounding between asymmetric users, and voice agents are criticised precisely for under-grounding. A surface that cannot ground and cannot graduate its own authority is not a collaboration surface; it is a remote control.
\end{keyinsight}

One structural gap sits underneath all four surfaces: the standards bodies lag the practice they would govern. The EU AI Act's oversight duties, NIST's risk-management framework, and ISO/IEC~42001 all extend pre-agentic control regimes rather than issue agent-native disclosure norms \parencite{EUAIAct2025, NISTAIRMF, ISO42001}, and no ratified W3C standard yet governs how an agent identifies itself or its authorising principal --- the identity and provenance work remains draft \parencite{AIIdentity2026, Souza2025ProvAgent}. That vacuum is why disclosure on the forum surface is enforced by after-the-fact bans rather than protocol, and why an ecosystem that binds every actor to one signing identity by construction is making a wager on where the standards will land, rather than following a settled one. The wager is worth stating plainly because it is falsifiable: if agent-identity standards ratify a different primitive, the canon's single-key model becomes a migration cost rather than a head start.

The four-surface taxonomy frames the component chapters that open with Chapter~\ref{ch:visionflow-canon}, where each DreamLab component is examined against the surface it serves, and it structures the gap register that follows in Chapter~\ref{ch:gap-register}. The register applies these four lenses systematically to the codebase, and the pattern the surfaces already predict holds there: the asynchronous decision surface is the most fully built and the most honestly audited, the visual surfaces render awareness without yet closing the steering loop, and voice governance is designed ahead of its implementation. Judgment must be legible, auditable, and deliberate, which is why the weight sits where it sits.
